\chapter*{ЗАКЛЮЧЕНИЕ}
\addcontentsline{toc}{chapter}{Заключение}

В ходе выполнения выпускной квалификационной работы разработана корпоративная
информационная система картографии и маршрутов Guide Helper. Практический
результат работы
представляет собой веб-приложение и набор серверных сервисов, которые позволяют
автору маршрута создать точки на карте, прикрепить к ним фотографии и заметки,
рассчитать характеристики маршрута, опубликовать ссылку и получить обратную
связь от пользователей.

Все десять задач, поставленных во введении, выполнены:
\begin{enumerate}
  \item проведён анализ предметной области цифрового туризма, выявлены ключевые потребности целевой аудитории~--- гидов, инструкторов и туристических организаций;
  \item выполнен сравнительный анализ пяти аналогов (Wikiloc, Google My Maps, Strava, Komoot, AllTrails), обоснована необходимость собственной разработки;
  \item сформулированы 37 функциональных и 6 нефункциональных требований к системе; все 37 функциональных требований реализованы (100\%);
  \item спроектирована архитектура системы с диаграммами контекста, контейнеров и ER-схемой базы данных; обоснован выбор Rust, Go, TypeScript, PostgreSQL, NATS JetStream и MinIO с учётом характера нагрузки отдельных компонентов;
  \item реализованы сервисы аутентификации и маршрутов на Rust/Axum с JWT-аутентификацией, ролевой моделью и полным циклом операций создания, чтения, изменения и удаления маршрутов с социальными функциями;
  \item разработано клиентское веб-приложение с интерактивной картой Leaflet, профилем высот, прогнозом погоды, каталогом маршрутов и административной панелью; подготовлена настольная версия на Tauri;
  \item реализован асинхронный конвейер обработки фотографий через NATS JetStream с хранением в MinIO, уведомлениями через WebSocket;
  \item интегрирован ИИ-ассистент с поддержкой вызова инструментов и несколькими провайдерами языковых моделей;
  \item настроен конвейер автоматической сборки, проверки и развёртывания через GitHub Actions с публикацией Docker-образов в GHCR и развёртыванием в Kubernetes;
  \item написано 115 модульных тестов, проведено функциональное тестирование по критическим путям трёх ролей, выполнен анализ производительности.
\end{enumerate}

Ключевое отличие системы Guide Helper по сравнению с рассмотренными аналогами
состоит не в отдельной функции, а в объединении нескольких действий в одном
сценарии: автор маршрута работает с точками, фотографиями, заметками,
характеристиками маршрута и публикацией в одном интерфейсе. Серверная обработка
изображений и ИИ-функции дополняют основной сценарий, но не заменяют его:
главным результатом остаётся публикуемый интерактивный маршрут.

Практическая значимость работы заключается в том, что гид или автор маршрута
может в одном интерфейсе подготовить интерактивный маршрут, прикрепить к точкам
фотографии и заметки, опубликовать ссылку и получить обратную связь от
пользователей. Для туриста система предоставляет возможность заранее оценить
маршрут по карте, фотографиям, описанию, сложности и пользовательским оценкам.
Развёртывание в Kubernetes-кластере и конвейер автоматической сборки относятся к
инженерной части результата: они позволяют воспроизводимо обновлять систему,
изолировать отдельные сервисы и сопровождать демонстрационный стенд.

Направлениями дальнейшего развития системы могут служить: развитие мобильного
клиента до стабильного решения с устойчивой фоновой записью маршрута и
офлайн-синхронизацией, экспорт маршрутов в PDF, персонализированные рекомендации
маршрутов и расширение инструментов ИИ-ассистента.
